\sectioncentered*{Заключение}
\addcontentsline{toc}{section}{Заключение}

Для предметной области музыки описан сценарий из восьми сообщений, заданы классы намерений \texttt{about\_entity}, \texttt{about\_performer}, \texttt{about\_author}, \texttt{about\_album} и подключены сущности БЗ. Классификация выполняется локальным сервисом, совместимым с API Wit.ai. В диалоговом окне Ники проверены вопросы <<Что такое~\ldots?>> и четыре реплики сценария.

Ответы на вопросы об исполнителе, авторе и альбоме зависят от состояния БЗ и задаются логическими правилами лабораторной работы №3.

\begin{thebibliography}{9}
	\addcontentsline{toc}{section}{Список использованных источников}

	\bibitem{nika}
	NIKA: Intelligent Knowledge-driven Assistant [Электронный ресурс]. \mdash  Режим доступа: \url{https://github.com/ostis-apps/nika}. \mdash  Дата доступа: 01.10.2026.

	\bibitem{wit}
	Wit.ai documentation [Электронный ресурс]. \mdash  Режим доступа: \url{https://wit.ai/docs}. \mdash  Дата доступа: 01.10.2026.

	\bibitem{wit2851}
	Wit.ai intent classification broken for newly created apps [Электронный ресурс]. \mdash  Режим доступа: \url{https://github.com/wit-ai/wit/issues/2851}. \mdash  Дата доступа: 01.10.2026.

	\bibitem{ostis}
	OSTIS: Open Semantic Technology for Intelligent Systems [Электронный ресурс]. \mdash  Режим доступа: \url{https://ostis.net/}. \mdash  Дата доступа: 01.10.2026.
\end{thebibliography}
